\chapter{Fundamente teoretice}
\label{ch:fundamente}

\section{Ierarhia de memorie și localitatea referințelor}

Ideea de cache provine din arhitectura calculatoarelor: o memorie mică și rapidă păstrează o
submulțime a datelor dintr-o memorie mare și lentă \cite{hennessy}. Eficiența ei se bazează pe
\emph{localitatea referințelor}: datele accesate recent (localitate temporală) sau populare vor fi
probabil accesate din nou. În aplicațiile web, ierarhia arată astfel:

\begin{table}[H]
\centering\small
\begin{tabular}{lll}
\toprule
Nivel & Exemplu & Latență tipică \\
\midrule
Memorie în proces (L1) & \texttt{HashMap} protejat de mutex & 0{,}05–1 µs \\
Cache în rețea locală (L2) & Memcached, Redis, \texttt{cache-node} & 50–300 µs \\
Bază de date, date în RAM & PostgreSQL, \emph{shared buffers} & 0{,}2–2 ms \\
Bază de date, agregări / disc & join + \texttt{GROUP BY}, I/O & 2–100 ms \\
\bottomrule
\end{tabular}
\caption{Ordine de mărime ale latențelor în ierarhia de date a unei aplicații web.}
\end{table}

Timpul mediu de acces efectiv este dat de relația clasică:
\begin{equation}
  T_{\text{mediu}} = h \cdot T_{\text{cache}} + (1-h)\cdot\left(T_{\text{cache}} + T_{\text{DB}}\right)
  = T_{\text{cache}} + (1-h)\cdot T_{\text{DB}},
  \label{eq:amat}
\end{equation}
unde $h$ este rata de hit. Deoarece $T_{\text{DB}} \gg T_{\text{cache}}$, câștigul este dominat de
$(1-h)$: trecerea de la $h=0{,}90$ la $h=0{,}99$ reduce de zece ori numărul de accesări ale bazei de
date. Rata de hit este deci metrica centrală a oricărui cache.

\section{Distribuția Zipf a popularității}

Studiile asupra traficului web \cite{breslau} arată că popularitatea obiectelor urmează aproximativ o
lege Zipf: probabilitatea ca o cerere să vizeze obiectul de rang $i$ dintr-un total de $N$ este
\begin{equation}
  p_i = \frac{i^{-s}}{H_{N,s}}, \qquad H_{N,s} = \sum_{k=1}^{N} k^{-s},
\end{equation}
unde exponentul $s$ (de obicei între $0{,}6$ și $1{,}2$) controlează asimetria. În modelul
\emph{Independent Reference Model} (IRM), cel mai bun cache static de capacitate $C$ păstrează cele
mai populare $C$ obiecte și atinge rata de hit
\begin{equation}
  h^{*}(C) = \frac{H_{C,s}}{H_{N,s}},
  \label{eq:ideal}
\end{equation}
care este o limită superioară utilă pentru politicile reale (folosită în \cref{sec:sim}). Pentru
$s = 0{,}99$ și $N = 100.000$, un cache cu doar 10\% din chei poate servi teoretic aproape 80\% din
cereri — de aici eficiența cache-urilor mici.

\section{Politici de evacuare}

Când cache-ul este plin, trebuie ales ce intrare se elimină. Algoritmul optim (Belady, MIN)
\cite{belady} elimină intrarea care va fi folosită cel mai târziu în viitor, dar necesită cunoașterea
viitorului. Politicile practice aproximează acest comportament:

\begin{description}
  \item[FIFO] elimină cea mai veche intrare inserată; ignoră complet accesele. Cost minim, rată de hit
        mai slabă.
  \item[LRU] (\emph{Least Recently Used}) elimină intrarea neaccesată de cel mai mult timp. Exploatează
        localitatea temporală, dar fiecare acces modifică structura (mutare în capul listei), ceea ce
        crește contenția sub concurență. Este vulnerabilă la \emph{scanări}: o parcurgere secvențială a
        multor chei reci elimină setul fierbinte.
  \item[CLOCK] (\emph{second chance}) \cite{corbato} aproximează LRU cu un singur bit de referință per
        intrare: la acces doar se setează bitul; la evacuare, intrările cu bitul setat primesc o „a doua
        șansă" (bitul se resetează). Accesul nu mai mută noduri, deci operația de citire este mai ieftină.
\end{description}

Politici mai sofisticate (LFU, ARC \cite{megiddo}, W-TinyLFU) combină recența cu frecvența; ele sunt
discutate ca direcții viitoare.

\section{Strategii de citire și scriere}

\paragraph{Cache-aside (lazy loading).} Aplicația caută întâi în cache; la \emph{miss} citește din baza
de date și populează cache-ul. Este modelul folosit pentru citiri în acest proiect.

\paragraph{Invalidare la scriere.} La o modificare, aplicația scrie în baza de date și apoi
\emph{șterge} intrarea din cache. Este simplă și sigură, dar următoarea citire va fi un miss.

\paragraph{Write-through.} Aplicația scrie în baza de date și apoi \emph{actualizează} cache-ul cu
valoarea nouă. Rata de hit se păstrează, cu prețul unei latențe de scriere mai mari.

\paragraph{Write-behind (write-back).} Aplicația actualizează imediat cache-ul și confirmă clientul;
scrierea în baza de date este amânată și grupată în loturi. Scrierile devin foarte rapide și numărul de
tranzacții scade drastic, dar consistența devine \emph{eventuală}, iar o cădere înainte de flush poate
pierde date.

\section{Consistență și coerență}

Un cache introduce inevitabil o a doua copie a datelor. Teorema CAP \cite{brewer} formalizează
compromisul dintre consistență și disponibilitate în prezența partițiilor de rețea; pentru cache-uri,
alegerea uzuală este \emph{consistența eventuală} cu fereastră de inconsistență limitată de TTL.
Probleme clasice:

\begin{itemize}
  \item \textbf{coerența cache-urilor locale}: fiecare instanță are propria copie; o scriere pe instanța A
        trebuie să invalideze copia din B. Proiectul folosește notificări \texttt{LISTEN/NOTIFY} emise de
        un trigger PostgreSQL.
  \item \textbf{cursa citire–invalidare}: un cititor obține valoarea veche din DB, un scriitor modifică și
        invalidează, apoi cititorul populează cache-ul cu valoarea veche. TTL-ul limitează durata
        inconsistenței; soluții complete folosesc \emph{lease}-uri \cite{memcache-fb} sau versiuni.
  \item \textbf{cache stampede} (\emph{thundering herd}): la expirarea unei chei foarte populare, sute de
        cereri concurente ratează simultan și lovesc baza de date. Soluția implementată este
        \emph{single-flight}: o singură cerere per cheie interoghează baza de date, celelalte așteaptă
        rezultatul.
\end{itemize}

\section{Distribuirea datelor: consistent hashing}

Distribuirea cheilor pe $n$ noduri prin $\text{hash}(k) \bmod n$ este echilibrată, dar la adăugarea unui
nod aproape toate cheile își schimbă nodul, ceea ce golește efectiv cache-ul. \emph{Consistent hashing}
\cite{karger} plasează nodurile și cheile pe un inel; fiecare cheie aparține primului nod întâlnit în
sens orar. La modificarea topologiei se mută doar $\approx 1/n$ din chei. Pentru echilibrare, fiecare nod
fizic apare de mai multe ori pe inel (\emph{noduri virtuale}).

\section{Metrici și legi de performanță}

\paragraph{Legea lui Little.} Pentru un sistem stabil, numărul mediu de cereri în sistem este
$L = \lambda W$ \cite{little}. Într-un test \emph{closed-loop} cu $L$ clienți concurenți, debitul este
$\lambda = L / W$: dacă latența rămâne constantă, debitul crește liniar cu concurența, până la saturarea
unei resurse.

\paragraph{Legea lui Amdahl.} Dacă o fracțiune $f$ a lucrului este paralelizabilă, accelerarea maximă pe
$p$ procesoare este \cite{amdahl}
\begin{equation}
  S(p) = \frac{1}{(1-f) + f/p}.
  \label{eq:amdahl}
\end{equation}
Secțiunile critice (de exemplu un mutex global) reprezintă partea serială $(1-f)$; sharding-ul reduce
această parte.

\paragraph{Percentile de latență.} Media ascunde coada distribuției. Pentru aplicații web contează p95,
p99 și p99.9: cu 100 de apeluri per pagină, aproape fiecare pagină „vede" latența p99. Latențele sunt
înregistrate cu histograme HDR \cite{tene}, cu precizie de 3 cifre semnificative.
